\chapter{函数的连续性}
\section{连续性概念}
\begin{definition}
    设函数$f(x)$在$U(x_0)$上有定义.若
    \begin{equation}
        \lim_{x\to x_0} f(x)=f(x_0),
    \end{equation}
    则称$f(x)$在点$x_0$连续.
\end{definition}
$\varepsilon - \delta$语言:$\forall \varepsilon>0,\exists \delta>0,\forall x\in U(x_0,\delta),|f(x)-f(x_0)|<\varepsilon.$

记$\Delta x=x-x_0,\Delta y=y-y_0$,则函数$f(x)$在点$x_0$连续等价于\[\lim_{\Delta x\to 0} \Delta y=0.\]

由函数$f(x)$在点$x=x_0$连续的定义可知,其必要条件是$f(x)$在$x\to x_0$时极限存在.
\begin{definition}
    设函数$f(x)$在$U(x_0)$上有定义.若函数$f(x)$在点$x_0$无定义,或$f$在点$x_0$有定义而不连续,则称点$x_0$为函数$f$的间断点.
\end{definition}
\begin{enumerate}
    \item 可去间断点:$f(x)$在$x\to x_0$时极限存在,但$f$在点$x_0$无定义,或有定义不等于极限值.
    \item 跳跃间断点:$f(x)$在点$x_0$的左右极限均存在,但不相等.
    \item 第二类间断点:$f(x)$在点$x_0$至少有一侧极限不存在.
\end{enumerate}

若函数在区间$I$上的每一点都连续(区间端点考虑左(右)极限),则称$f$为区间$I$上的连续函数.

\begin{example}
    证明:黎曼函数
    \begin{equation*}
        R(x)=\begin{cases}
            \dfrac{1}{q},\text{当}x=\dfrac{q}{p}(p,q\text{为正整数},\dfrac{q}{p}\text{为既约真分数})\\
            0,\text{当}x=0,1\text{及}(0,1)\text{上的无理数}
        \end{cases}
    \end{equation*}
    在$0,1$及$(0,1)$上的无理数点都连续,$(0,1)$上的有理数点都不连续.
\end{example}
\begin{proof}
    任取$x_0\in(0,1)$,$\forall \varepsilon>0$,满足$\dfrac{1}{q}\geq\varepsilon$,即$q\leq\dfrac{1}{\varepsilon}$的$q$只有有限个,从而$\dfrac{q}{p}$也只有有限个,设为$x_1,x_2,\cdots,x_n$.\\
    取$\delta_1=\min\{|x_0-x_i|:i=1,2,\cdots,n\text{且}x_0\neq x_i\}$,$\delta=\min\{x_0,1-x_0,\delta_1\}$,当$x\in U^{\circ}(x_0,\delta)$时,$|R(x)-0|<\varepsilon$.
    因此$\forall x_0\in(0,1)\text{有}\displaystyle\lim_{x\to x_0}R(x)=0$.于是可知对于$(0,1)$上的无理点,函数极限值均等于函数值$0$,连续;而对于$(0,1)$上的有理点,函数极限值一定不会等于函数值$\dfrac{1}{q}$,不连续.\\
    当$x_0=0$时,$\forall \varepsilon>0$,满足$\dfrac{1}{q}\geq\varepsilon$,即$q\leq\dfrac{1}{\varepsilon}$的$q$只有有限个,从而$\dfrac{q}{p}$也只有有限个,设为$x_1,x_2,\cdots,x_n$.取$\delta=\min\{x_i:i=1,2,\cdots,n\text{且}x_i\neq 0\}$,当$x\in U^{\circ}_{+}(x_0,\delta)$时,$|R(x)-0|<\varepsilon$.因此$\displaystyle\lim_{x\to 0}R(x)=0=R(0)$,连续.当$x_0=1$时同理.
\end{proof}

\section{连续函数的性质}
\begin{theorem}[局部有界性]
    若函数$f$在点$x_0$连续,则$f$在某$U(x_0)$上有界.
\end{theorem}
\begin{theorem}[局部保号性]
    若函数$f$在点$x_0$连续,且$f(x_0)>0$,则对任何$f(x_0)>r$,存在某$U(x_0)$,使得对一切$x\in U(x_0)$,有$f(x)>r$.
\end{theorem}
\begin{theorem}
    若函数$f$和$g$在点$x_0$连续,则$f\pm g,f\cdot g,f/g$也都在点$x_0$连续.
\end{theorem}
\begin{theorem}[有界性定理]
    若函数$f$在闭区间$[a,b]$上连续,则$f$在闭区间$[a,b]$上有界.
\end{theorem}
\begin{proof}
    反证:假设$f(x)$在$[a,b]$上无界,那么存在$x_n\subset[a,b]$,使得
    \begin{equation*}
        f(x_n)>n,n=1,2,\cdots
    \end{equation*}
    于是有$\displaystyle\lim_{n\to\infty}f(x_n)=+\infty$.\\
    另一方面,由$x_n\subset[a,b]$有界,根据致密性定理,$\{x_n\}$有收敛子列$\{x_{n_k}\}$,设$\displaystyle\lim_{k\to\infty}x_{n_k}=x_0$.
    \\由$a\leq x_{n_k}\leq b$可知$a\leq x_0\leq b$,故$f(x)$在点$x_0$连续.根据归结原则可得
    \begin{equation*}
        +\infty=\displaystyle\lim_{n\to\infty}f(x_n)=\lim_{k\to\infty}f(x_{n_k})=\lim_{x\to x_0}f(x)=f(x_0),
    \end{equation*}
    矛盾.
\end{proof}
\begin{theorem}[最大,最小值定理]
    若函数$f$在闭区间$[a,b]$上连续,则$f$在闭区间$[a,b]$上有最大值和最小值.
\end{theorem}
\begin{proof}
    由有界性定理,确界原理,$f$在$[a,b]$上存在上确界,$\displaystyle\sup_{x\in[a,b]}f(x)=M$.下证:存在$x_0\in[a,b]$使得$f(x_0)=M$.\\
    $\forall x\in[a,b],f(x)\leq M$,且任取$\varepsilon_n=\dfrac{1}{n}(n=1,2,\cdots)$,则$\exists x_n\in[a,b]$使得$f(x_n)>M-\varepsilon_n\Rightarrow M-\dfrac{1}{n}<f(x_n)\leq M\Rightarrow\displaystyle\lim_{n\to\infty}f(x_n)=M$.$\{x_n\}$有界,根据致密性定理,$\{x_n\}$有收敛子列$\{x_{n_k}\}$,设$\displaystyle\lim_{k\to\infty}x_{n_k}=x_0$.
    \begin{equation*}
        \displaystyle M=\lim_{n\to\infty}f(x_n)=\lim_{k\to\infty}f(x_{n_k})=\lim_{x\to x_0}f(x)=f(x_0).
    \end{equation*}
    同理可证$f$在$[a,b]$上存在最小值.
\end{proof}
\begin{theorem}[介值性定理]
    设函数$f$在闭区间$[a,b]$上连续,且$f(a)\neq f(b)$.若$\mu$为介于$f(a)$与$f(b)$之间的任何实数,则至少存在一点$x_0\in(a,b)$,使得$f(x_0)=0$.
\end{theorem}
\begin{proof}
    不妨设$f(a)<\mu<f(b)$.令$g(x)=f(x)-\mu$,则$g$也是$[a,b]$上的连续函数,且$g(a)<0,g(b)>0$,于是定理的结论转化为:存在$x_0\in(a,b)$,使得$g(x_0)=0$.设集合
    \begin{equation*}
        E=\{x|g(x)<0,x\in[a,b]\}.
    \end{equation*}
    显然$E$为非空有界数集,根据确界原理,$E$存在上确界$x_0=\sup E$.\\
    由$g(a)<0$,根据连续函数的局部保号性,$\exists \delta>0,\forall x\in U^{\circ}_{+}(a,\delta),g(x)<0$,于是$\displaystyle g(a+\dfrac{\delta}{2})<0$,即$\displaystyle a+\dfrac{\delta}{2}\in E$,而$\displaystyle a<a+\dfrac{\delta}{2}\leq \sup E=x_0$,因此$x_0\neq a$;由$g(b)>0$,根据连续函数的局部保号性,$\exists \delta>0,\forall x\in U^{\circ}_{-}(b,\delta),g(x)>0$,于是$\forall x\in E$,一定有$x\in[a,b-\delta]$,而$b>b-\delta\geq\sup E=x_0$,因此$x_0\neq b$.这说明$x_0\in(a,b)$,下证$g(x_0)=0$.\\
    反证:若$g(x_0)\neq 0$,则$g(x_0)>0$或$g(x_0)<0$.\\
    若$g(x_0)<0$,则$\exists\eta>0,\forall x\in(x_0,x_0+\eta)\subseteq(a,b),g(x)<0$,从而有$\displaystyle g(x_0+\dfrac{\eta}{2})<0$,即$\displaystyle x_0+\dfrac{\eta}{2}\in E$.矛盾;若$g(x_0)>0$,则存在$\eta>0,\forall x\in(x_0-\eta,x_0]\subseteq(a,b),g(x)>0$,另一方面,由$x_0=\sup E$可知$\forall x\in(x_0,b],x>x_0$,即$x\notin E,g(x)\geq 0$.综合可得,$\forall x\in(x_0-\eta,b],g(x)\geq 0$,于是$E\subseteq[a,x_0-\eta]$,说明$x_0=\sup E\leq x_0-\eta$.矛盾.因此$g(x_0)=0$.
\end{proof}
\begin{theorem}[根的存在定理]
    若函数$f$在闭区间$[a,b]$上连续,且$f(a)$与$f(b)$异号,则至少存在一点$x_0\in(a,b)$,使得$f(x_0)=0$.
\end{theorem}
\begin{definition}
    设$f$为区间$I$上的函数.若对任给的$\varepsilon>0$,存在$\delta=\delta(\varepsilon)>0$,使得对任何$x',x''\in I$,只要$|x'-x''|<\delta$,就有
    \begin{equation*}
        |f(x')-f(x'')|<\varepsilon,
    \end{equation*}
    则称函数$f$在区间$I$上一致连续.
\end{definition}
上述定义说明如果$f$在区间$I$上一致连续,那么$f$在区间$D\subseteq I$上也一致连续.
\begin{theorem}
    函数$f(x)$定义在区间$I$上.$f(x)$在$I$上一致连续的充要条件为:对任何数列$\{x_n\},\{y_n\}\subset I$,若$\displaystyle\lim_{n\to\infty}(x_n-y_n)=0$,则
    \begin{equation*}
        \displaystyle\lim_{n\to\infty}[f(x_n)-f(y_n)]=0.
    \end{equation*}
\end{theorem}
\begin{proof}
    必要性:若$f(x)$在$I$上一致连续,则$\forall\varepsilon>0,\exists\delta(\varepsilon)>0,\forall x',x''\in I,|x'-x''|<\delta$,则$|f(x')-f(x'')|<\varepsilon$.任取$\{x_n\},\{y_n\}\subset I$,满足$\displaystyle\lim_{n\to\infty}(x_n-y_n)=0$,那么对于上述的$\delta>0,\exists N>0,\forall n>N,|x_n-y_n|<\delta$,于是有$|f(x_n)-f(y_n)|<\varepsilon$,即$\displaystyle\lim_{n\to\infty}[f(x_n)-f(y_n)]=0$.\\
    充分性:反证:若$f(x)$在$I$上不一致连续,则$\exists \varepsilon_0>0,\forall\delta>0,\exists x',x''$,满足$|x'-x''|<\delta$,但$|f(x')-f(x'')|\geq\varepsilon_0$.取$\displaystyle\delta_n=\dfrac{1}{n}$,则$\displaystyle\exists\{x_n\},\{y_n\}\subset I,|x_n-y_n|<\delta=\dfrac{1}{n}$,有$|f(x_n)-f(y_n)|\geq\varepsilon_0$,于是$\displaystyle\lim_{n\to\infty}(x_n-y_n)=0$,但$\displaystyle\lim_{n\to\infty}[f(x_n)-f(y_n)]\neq 0$.矛盾,所以$f(x)$在$I$上一致连续.
\end{proof}
\begin{theorem}[康托定理]
    若函数$f$在$[a,b]$连续,则$f$在$[a,b]$一致连续.
\end{theorem}
\begin{theorem}
    设区间$I_1$的右端点为$c\in I_1$,区间$I_2$的左端点为$c\in I_2$.若$f$分别在$I_1$和$I_2$上一致连续,则$f$在$I=I_1\cup I_2$上也一致连续.
\end{theorem}
\begin{proof}
    由$f$在$I_1$连续可知,$\forall\varepsilon>0,\exists\delta_1>0,\forall x',x''\in I_1$,只要$|x'-x''|<\delta_1,$就有$|f(x')-f(x'')|<\varepsilon$;
    由$f$在$I_2$连续可知,$\forall\varepsilon>0,\exists\delta_2>0,\forall x',x''\in I_2$,只要$|x'-x''|<\delta_2,$就有$|f(x')-f(x'')|<\varepsilon$.另一方面,$f$在$I_1$和$I_2$一致连续,可知$f$在$x=c$点连续,即$\displaystyle\exists\delta_3>0,\forall x\in U(c,\delta_3),|f(x)-f(c)|<\dfrac{\varepsilon}{2}$.\\
    取$\delta=\min\{\delta_1,\delta_2,\delta_3\},\forall x',x''\in I$,且$|x'-x''|<\delta$,
    \begin{enumerate}
        \item 若$x',x''$同属于$I_1$或$I_2$,则有$|f(x')-f(x'')|<\varepsilon$
        \item 若$x'\in I_1,x''\in I_2$,此时有$|x'-c|\leq |x'-x''|<\delta_3$,$|x''-c|\leq |x'-x''|<\delta_3$,于是$\displaystyle|f(x')-f(x'')|\leq |f(x')-f(c)|+|f(x'')-f(c)|<\dfrac{\varepsilon}{2}+\dfrac{\varepsilon}{2}=\varepsilon$
    \end{enumerate}
    因此$f$在$I=I_1\cup I_2$上也一致连续.
\end{proof}
